\documentclass[12pt]{article}

\usepackage[a4paper, total={6in, 8in}]{geometry}
\usepackage{textcomp}

\begin{document}
\setlength{\parindent}{0pt}

\def \t {\textrightarrow}
\def \v {\vspace{1em}}
\def \bi {\begin{itemize}}
\def \ei {\end{itemize}}
\def \s[#1] {\section*{#1}}
\def \ss[#1] {\subsection*{#1}}
\def \sss[#1] {\subsubsection*{#1}}

\s[Recuperare 16 marzo]
\s[D'annunzio - Il piacere]
D'annunzio \t l'inettitudine, non ha la forza di reagire \t il super uomo \\
Forte concentrazione di lessico che stordisce il lettore quasi \t si viene rapiti dalla parola, il verso è tutto \t che rende inimitabile l'opera di d'annunzio \\
Il personaggio del fuoco matura la sua decisione, riflette \t e sulla base di un attenta riflessione e con un analisi del doppio, il personaggio Laurispa ripercorre la sua vita e viive un passato diverso dal suo \t abbraccia Ippolita, e uccide lei e se stesso \\
Amore come ossessione \t dietro c'è l esperienza di ?? = "Le cult du moi" \t la lettura sul pan faltal di Tristano e Isotta che lui compie \\
Il mito d'annunziano dell'uomo invincibile frana \t questo romanzo avvia la parabola dell'inetto

\v

Nel nostro 900 il romanzo tende a portare la psicanalisi come filo conduttore e campo di esistenza \t il trionfo della morte, e il lavoro che caratterizza l'autonalisi del personaggio (Laurispia) \t e i timori a comprende da parte di Ippolita, ovvero la donna amata da lui \\
La psicologia profonda dei personaggi si lascia permeare da un gusto europeo \t si sfronda la tradizione, si analizza la duttilita dello stato d animo \\
Incipit con tensione ancora verista \t frana \\
Trionfo della morte, dove si ha un ribaltamento del romanzo tradizionale + Dostojevski \\
Parlare di una morale nei suoi romanzi è fallace \t l'autonalisi, l'inetto \t e si smonta la fan fatal, che caratterizza d'annunzio \\
Con il trionfo della morte, Ippoolita si presenta come la nemica \t Laurispa è ispirato da Nietzsche, me anche limitato da lui \t vuole andare oltre

\ss[Le vergini delle roccie]
Parla di Caludio Cantello, che è un artista \t ma disprezza e disdegna la volgarita borghese, e la bassezza dei costumi, e questo mondo cosi arroccato alle tradizioni popolari \\
Cantello rappresenta un altro eroe estetizzante per l'arte \t è uno scultore, ed è intenzionato in senso antidemocratico ad affermare una nuova personalita \\
Pensa a un processo di selezione e combinazione con una ipoetitica figura femminile, che avrebbe fatto procreare il futuro superuomo \\
Questo inugura i "Romanzi del giglio" \t qua l'attenzione si sposta verso la presa di consapevolezza di una fiducia da guadnare \t menre "Romanzi della rosa" focus sulla passione \\
Questo è il romanzo che mostra la crudelta prevaricante del super uomo \t l'opera d'arte è uno studio di come l'uomo deve essere procreato \\
La creazione di un essere perfetto, che possa stare a capo di una nuova roma \\
Cantello raccoglie gli stralci del Perelli e dell Aurispa \t ed è inetto, perche cerca la perfezione che pero non esiste in terra \\
Spirito prevaricante e il fatto che tutti i valori debbano essere sconfessati \\
Egotismo \t d'annunzio pone al centro il suo ego in ogni atteggiamento \t che esclude l'empatia, e mettersi nei panni dell'altro

\ss[Il fuoco]
La ricerca di d'annunzio non si ferma \t e il fuoco \t la passione che brucia \\
Stelio Effrena \t protagonista del fuoco, che esce nell anno 1900 \\
Ha evidenti tratti superomistici \t superuomo ed esteta si fondono \\
L'attivita di d'annunzio come romanziere si ripropone \t d'annunzio si fa propaganda \\
Protagonista è un giovane brillante, intellettuale \t che vuole creare una opera d'arte totalizzante, che unisca musica, lett. etc. \t e il suo obbiettivo è quello di sconfiggere la morte raggiungendo la gloria \\
Proprio qui c'è una condesazione di poetica wagneriana \t anche i musicisti hanno una loro poetica \t poetica wagneriana è un fuoco di passione che caratterizza il pianeta marte \\
Incarna con l'amore per wagner un idolo del tempo \t wagner e l'Effrena, erano fortemente legati \t e proprio qui si descrive la partecipazione al funerale di Wagner, e l'effettiva ripdozuione della venenzia in cui Stelio Effrena che porta la bara di Wagner = conniugo intensissimo \\
La donna si chiama Foscarina (che è una proiezione di Aduse) \t il collegamento è che lei fa l'attrice, e all apice della carriera si trova divisa tra il successo personale e quello della carriera \\
Effrena quindi è invece la proiezione di d'annunzio stesso \\
L'aduse in qualche modo scalpita \t e la relazione tra lui e l'aduse non viene annunciata platealmente \t ma chi seguiva l'autore vede l'anticipazione della fine del rapporto in questo romanzo
\end{document}
